\BourbakiExercises{习题}

\begin{enumerate}
\def\labelenumi{\arabic{enumi})}
\item
  设 A 是环，m 是 A 的一个极大左理想。设 B 是 A 中使 \(mb \subset m\) 的元素 b 组成的集合。证明 B 是 A 的包含 m 的、且 m 在其中是双边理想的最大子环。设 S 是 A-模 \(A_{s}/m\)，D 是它的交换子，\(\pi\) 是从 \(A_{s}\) 到 S 的典范映射。证明：给定 \(b \in B\)，存在唯一元素 \(u_{b}\) 属于 \(\operatorname{End}_{\mathrm{A}}(\mathrm{S})\)，使 \(u_{b}(\pi(a)) = \pi(ab)\) 对每个 \(a \in A\) 成立。映射 \(b \mapsto u_{b}\) 诱导出从 B/m 到 \(\operatorname{End}_{\mathrm{A}}(\mathrm{S})\) 的一个同构。
\item
  设 D 是体，V 是一个右 D-向量空间，A 是 \(\operatorname{End}_{\mathrm{D}}(\mathrm{V})\) 的一个包含有限秩自同态的子环（例如 \(A = \operatorname{End}_{\mathrm{D}}(\mathrm{V})\)）。证明对偶 \(V^{*} = \operatorname{Hom}_{\mathrm{D}}(\mathrm{V}, \mathrm{D})\) 作为右 A-模是单模。
\item
  设 D 是体，V 是一个右 D-向量空间，A 是 V 的自同态环。对 V 中每条直线 L，用 \(m_{L}\) 表示 A 中使 \(u(L)=0\) 的元素 u 组成的集合。证明 \(m_{L}\) 是 A 的一个极大理想，我们有 \(m_{L}\neq m_{L'}\)（对 \(L\neq L'\)），但单 A-模 \(A/m_{L}\) 都同构于 A-模 V。
\item
  证明 Z-模 Q 既无单子模亦无单商模。
\item
  设 \(\mathrm{A}\) 是环。
\end{enumerate}

\begin{enumerate}
\def\labelenumi{\alph{enumi})}
\tightlist
\item
  设 M、N 是单生成的 A-模。证明下列性质等价：
\end{enumerate}

\begin{enumerate}
\def\labelenumi{(\roman{enumi})}
\item
  A-模 M 与 N 同构。
\item
  对每个生成元 \(m\)（\(\mathrm{M}\) 的），存在一个生成元 \(n\)（\(\mathrm{N}\) 的），它与 \(m\) 有相同的零化子。
\item
  存在 M 的一个生成元与 N 的一个生成元，它们有相同的零化子。
\end{enumerate}

\begin{enumerate}
\def\labelenumi{\alph{enumi})}
\setcounter{enumi}{1}
\item
  设 \(\mathfrak{a},\mathfrak{b}\) 是 A 的（左）理想。A-模 A/ \(\mathfrak{a}\) 与 A/ \(\mathfrak{b}\) 同构当且仅当存在 A 的一个元素 \(a\)，使 \(Aa + \mathfrak{b} = A\)，且 \(\mathfrak{a}\) 是 A 的元素 \(x\)（满足 \(xa\in \mathfrak{b}\)）组成的集合（利用 \(a\)）。
\item
  设 \(\mathfrak{m}\) 与 \(\mathfrak{n}\) 是 A 的极大左理想。单 A-模 A/ \(\mathfrak{m}\) 与 A/ \(\mathfrak{n}\) 同构当且仅当存在一个元素 \(a\)（属于 A- \(\mathfrak{n}\)），使 \(\mathfrak{ma} \subset \mathfrak{n}\) 。d) 由 b) 推出：A-模 \(\mathrm{A}_s\) 的同构于一个给定模的商 A-模组成的集合，其基数小于 \(\operatorname{Card}(\mathrm{A})\) 。
\end{enumerate}

\begin{enumerate}
\def\labelenumi{\arabic{enumi})}
\setcounter{enumi}{5}
\item
  设 M 是一个 A-模，使得对 M 中每个 \(x \neq 0\)，模 Ax 都是单模。证明：或者 M 是单模，或者位似环 \(A_{M}\) 是体（若 M 不是单模，考虑 M 的两个非零元素 x、y，使得 \(y \notin Ax\)，并考虑 \(x + y\) 的零化子）。
\item
  设 S 是一个单 A-模，其对偶 \(S^{*}\) 与双对偶 \(S^{**}\) 都是单模。证明映射 \(u \mapsto {}^{t}u\) 是从 \(\operatorname{End}_{\mathrm{A}}(\mathrm{S})\) 到 \(\operatorname{End}_{\mathrm{A}}(\mathrm{S}^{*})\) 的反体的一个同构。
\end{enumerate}

¶ 8) 设 A 是一个环，使得 A-模 \(A_{s}\) 与 \(A_{d}\) 都长度有限。我们把 \(A_{d}\) 与 \(A_{s}\) 的对偶等同起来。对 A 的每个左理想 s，s 的正交 \(s^{0}\) 是 s 在 A 中的右零化子；同样，对每个右理想 d，正交 \(d^{0}\) 是 d 的左零化子。证明下列性质等价：

\begin{enumerate}
\def\labelenumi{(\roman{enumi})}
\item
  每个（左或右）单 A-模都是自反的。
\item
  任何（左或右）单 A-模的对偶都是单模。
\item
  对每个极小左理想 \(\mathfrak{s}\) 与每个极小右理想 \(\mathfrak{d}\)，我们有 \(\mathfrak{s}^{00} = \mathfrak{s}\) 与 \(\mathfrak{d}^{00} = \mathfrak{d}\) 。
\item
  映射 \(\mathfrak{s} \mapsto \mathfrak{s}^0\) 是从 \(A\) 的左理想组成的集合到右理想组成的集合的一个双射。
\end{enumerate}

当这些性质成立时，我们称环 A 是自内射的（参见 X, §1, 第 171 页，习题 26）。

\begin{enumerate}
\def\labelenumi{\alph{enumi})}
\tightlist
\item
  设 A 是自内射的；设 M 是一个有限生成的（左或右）A-模。证明下列性质成立：
\end{enumerate}

(i') A-模 M 是自反的。

(ii') \(\mathbf{M}^*\) 的长度等于 \(\mathbf{M}\) 的长度。

(iii') 对每个子模 \(\mathrm{N}\)（属于 \(\mathrm{M}\)），我们有 \(\mathrm{N}^{00} = \mathrm{N}\) 。

(iv') 映射 \(N \mapsto N^{0}\) 是从 M 的子模组成的集合到 \(M^{*}\) 的子模组成的集合的一个双射。

\begin{enumerate}
\def\labelenumi{\arabic{enumi})}
\setcounter{enumi}{8}
\tightlist
\item
  单 A-模 S 的余长度是 S 的反模（在体 \(\operatorname{End}_{\mathrm{A}}(\mathrm{S})\) 上）的维数。我们称一个有限长度的 A-模 M 具有有限余长度，如果它有一个诸逐次商都有有限余长度的若尔当–赫尔德列；这些余长度之和称为 M 的余长度，记作 \(\lambda(\mathrm{M})\) 。若 N 是 M 的一个子模，则我们有 \(\lambda(\mathrm{M}) = \lambda(\mathrm{N}) + \lambda(\mathrm{M}/\mathrm{N})\) 。
\end{enumerate}

设 A 是一个自内射环（习题 8），使得模 \(A_{s}\) 与 \(A_{d}\) 都有有限余长度。若对每个单 A-模 S 我们都有 \(\lambda(S) = \lambda(S^{*})\)，则称 A 是弗罗贝尼乌斯环。

\begin{enumerate}
\def\labelenumi{\alph{enumi})}
\item
  设 M 是一个有限长度的 A-模。若 A 是弗罗贝尼乌斯环，则我们有 \(\lambda(M) = \lambda(M^{*})\) 。此外，若 A 是交换域 K 上一个有限秩代数，则我们有 \(\dim_K(M) = \dim_K(M^*)\)（利用习题 7）。
\item
  设 K 是交换域，A 是 \(\mathbf{M}_{6}(\mathrm{K})\) 的一个子代数，以元素 \(E_{11} + E_{55}\)、\(E_{12} + E_{56}\)、\(E_{21} + E_{65}\)、\(E_{22} + E_{66}\)、\(E_{33} + E_{44}\)、\(E_{13}\)、\(E_{23}\)、\(E_{45}\)、\(E_{46}\)（其中 \(E_{ij}\) 是 \(\mathbf{M}_{6}(\mathrm{K})\) 的矩阵单位）为基。证明 A 是自内射代数但不是弗罗贝尼乌斯代数（注意 A 的极小左理想是 \(KE_{13} + KE_{23}\)、\(KE_{45}\)、\(KE_{46}\)，而极小右理想是 \(KE_{13}\)、\(KE_{23}\)、\(KE_{45} + KE_{46}\)）。
\end{enumerate}

\begin{enumerate}
\def\labelenumi{\arabic{enumi})}
\setcounter{enumi}{9}
\tightlist
\item
  设 \(\rho : \mathrm{A} \to \mathrm{B}\) 是一个环同态，用它把 B 看作一个右 A-模。假设左 B-模 \(\operatorname{Hom}_{\mathrm{A}}(\mathrm{B}, \mathrm{A})\) 同构于 \(\mathrm{B}_s\) 。
\end{enumerate}

\begin{enumerate}
\def\labelenumi{\alph{enumi})}
\item
  设 M 是一个 B-模，\(\varphi : \mathrm{B}_s \to \mathrm{Hom}_{\mathrm{A}}(\mathrm{B}, \mathrm{A})\) 是一个同构。证明映射 \(u \mapsto \varphi(1) \circ u\)（从 \(\mathrm{Hom}_{\mathrm{B}}(\mathrm{M}, \mathrm{B})\) 到 \(\mathrm{Hom}_{\mathrm{A}}(\mathrm{M}, \mathrm{A})\)）是双射。
\item
  此外，设 A-模 B 是有限生成的。环 B 自内射（习题 8）当且仅当 A 也如此。
\item
  设 G 是一个有限群。证明：当把 B 取作赋有这样环结构的 A-模 A{[}G{]}——其中 \((ae_{g})(a^{\prime}e_{g^{\prime}})=aa^{\prime}e_{gg^{\prime}}\)（对 A 中 \(a,a^{\prime}\) 与 G 中 \(g,g^{\prime}\)）——时，上述条件满足。特别地，若 K 是交换域，则代数 K{[}G{]} 是自内射的。
\end{enumerate}

